% ====================================================================
\chapter{Implementación y Desarrollo}
\label{cap:implementacion}
% ====================================================================

\section{Implementación del subsistema de reconocimiento gestual}

\subsection{Preprocesamiento de datos y pipeline de extracción}
El pipeline procesa fotogramas de entrada mediante \textit{MediaPipe Holistic} para estructurar un vector instantáneo de 225 coordenadas ($33 \times 3$ de pose, $21 \times 3$ de mano izquierda y $21 \times 3$ de mano derecha). Para corregir discrepancias en la relación de aspecto de las cámaras ($4:3$ versus $16:9$), se aplica un reencuadre con bandas de relleno a una resolución uniforme de $853 \times 480$ píxeles previa extracción. Cada vector se centra en el punto medio de los hombros y se normaliza dividiendo por la distancia inter-hombros.

\subsection{Arquitectura neuronal TinySkeletonClassifier}
La clasificación de la secuencia temporal de 16 fotogramas se ejecuta mediante el modelo \textit{TinySkeletonClassifier}, cuya implementación comprende:
\begin{enumerate}
    \item Dos capas de convolución unidimensional temporal (núcleo de tamaño 3, 128 canales, normalización por lotes, activación ReLU y regularización Dropout de 0,2).
    \item Un módulo de codificación posicional sinusoidal sumado a la representación de 128 dimensiones.
    \item Un codificador Transformer de 2 capas con 4 cabezas de atención.
    \item Una capa de agregación por atención (\textit{Attention Pooling}) que asigna un escalar de relevancia $a_t$ a cada fotograma $t$, reduciendo la secuencia temporal a un vector único de 128 dimensiones:
    \begin{equation}
    \mathbf{v} = \sum_{t=1}^{T} a_t \mathbf{h}_t
    \label{eq:attention_pooling}
    \end{equation}
    \item Una capa densa lineal final hacia las 100 clases del vocabulario gestual.
\end{enumerate}

\section{Implementación del subsistema de traducción semántica bidireccional}

\subsection{Selección y especificación de la suite de modelos de lenguaje}
Para la traducción semántica bidireccional se adaptaron modelos de lenguaje pequeños de entre 0,5B y 3B parámetros:
\begin{itemize}
    \item \textbf{Familia Qwen2.5:} Variantes de instrucciones de 0,5B, 1,5B y 3B parámetros.
    \item \textbf{Familia Llama 3.2:} Variante de instrucciones de 1B parámetros.
    \item \textbf{Familia SmolLM2:} Variante de instrucciones de 1,7B parámetros.
\end{itemize}

\subsection{Régimen de entrenamiento y ajuste LoRA}
El ajuste fino se ejecutó dividiendo el corpus en 85\,\% para entrenamiento y 15\,\% para evaluación. La adaptación aplicó la técnica LoRA con rango $r = 16$ y alfa $\alpha = 32$ sobre las proyecciones de atención (\texttt{q\_proj}, \texttt{k\_proj}, \texttt{v\_proj}, \texttt{o\_proj}) y de la red feed-forward (\texttt{gate\_proj}, \texttt{up\_proj}, \texttt{down\_proj}). El entrenamiento utilizó el optimizador AdamW de 8 bits con tasa de aprendizaje de $3 \times 10^{-4}$ a $5 \times 10^{-4}$, programación de decaimiento coseno y precisión \textit{bfloat16}. La función de pérdida aplicó una máscara de gradiente para calcular el error exclusivamente sobre la respuesta esperada, ignorando los tokens del prompt de contexto.

\subsection{Canal directo (LSA $\rightarrow$ Español)}
El canal directo transforma secuencias de glosas en oraciones en español rioplatense. Las reglas incorporadas durante el ajuste incluyen la conversión de pronombres, la adición de marcas de género (\texttt{A} para femenino, \texttt{O} para masculino), el ajuste de tiempos verbales a partir de tokens temporales (\texttt{PASADO}, \texttt{FUTURO}) y el reordenamiento de la sintaxis \textit{Topic-Comment} hacia la estructura convencional del español.

\subsection{Canal inverso (Español $\rightarrow$ LSA)}
El canal inverso procesa la oración en español y emite la secuencia de glosas ordenada bajo la sintaxis \textit{Topic-Comment} (tiempo, lugar, sujeto, sustantivo, posesivo, numeral, verbo y partícula final de negación o pregunta). Los nombres propios y cifras se conservan en tokens únicos consolidados para su posterior procesamiento por la regla de deletreo dactilológico.

\subsection{Cuantización GGUF Q4\_K\_M y estimación de memoria}
Los modelos ajustados se fusionaron y convirtieron a binarios GGUF con cuantización \texttt{Q4\_K\_M}. La memoria RAM requerida durante la inferencia se calcula mediante la relación:
\begin{equation}
\text{RAM} \approx \mathbf{W} + KV_{\text{fp16}}(ctx=2048) + \text{Buffer}
\label{eq:ram_runtime}
\end{equation}
donde $\mathbf{W}$ representa el tamaño de los tensores cuantizados ($\sim 0{,}5$ bytes por parámetro), $KV_{\text{fp16}}$ la caché de claves y valores para 2048 tokens de contexto, y el buffer reserva espacio para las activaciones intermedias.

\section{Integración en Google Meet y extensión de navegador}

\subsection{Despliegue bajo Manifest V3}
La extensión web distribuye su lógica entre un script de contenido (\texttt{inject-gum.js}), un script en segundo plano (\textit{service worker}), un documento oculto (\textit{offscreen}) y un marco aislado (\textit{sandbox}). El esquema de paquetes y comunicación se ilustra en la Figura~\ref{fig:paquetes_extension}.

\begin{figure}[H]
    \centering
    \begin{tikzpicture}[
        font=\footnotesize,
        pkg/.style={draw, rounded corners=2pt, align=center, inner sep=4pt, fill=white, text width=3.3cm},
        arr/.style={-{Stealth[length=2.2mm]}, thick}
    ]
        \node[pkg] (sw) {\textit{Service worker}};
        \node[pkg, below left=0.7cm and 0.15cm of sw] (cs) {Script de contenido\\\texttt{inject-gum.js}};
        \node[pkg, below=0.7cm of cs] (canvas) {Lienzo\\\texttt{canvas}};
        \node[pkg, below right=0.7cm and 0.15cm of sw] (off) {Documento \textit{offscreen}};
        \node[pkg, below=0.7cm of off] (sb) {\textit{Sandbox}\\MediaPipe Holistic};
        \node[draw, rounded corners=3pt, inner sep=8pt, fit=(sw)(cs)(off)(sb)(canvas),
              label={[font=\scriptsize]above:Extensión Manifest V3}] (ext) {};
        \node[pkg, right=2.2cm of off] (ilsa) {Motor ILSA\\\texttt{127.0.0.1:8765}};
        \draw[arr, dashed] (sw) -- (cs);
        \draw[arr, dashed] (sw) -- (off);
        \draw[arr] (cs) -- node[left, font=\scriptsize] {cámara} (canvas);
        \draw[arr] (off.south) -- (sb.north);
        \draw[arr] (sb.west) to[out=180, in=180, looseness=1.15]
            node[left, font=\scriptsize] {225 puntos} (off.west);
        \draw[arr] (off.east) -- node[above, font=\scriptsize] {\texttt{POST /sign}} (ilsa.west);
        \coordinate (bajo) at ($(ilsa.south |- ext.south)+(0,-0.45)$);
        \coordinate (llega) at (bajo -| canvas.south);
        \draw[arr] (ilsa.south) -- (ilsa.south |- bajo) -- (llega) -- (canvas.south);
        \node[font=\scriptsize, below] at ($(bajo)!0.5!(llega)$) {clip};
    \end{tikzpicture}
    \caption{Paquetes de la extensión. La extracción corre en el \textit{sandbox}. El lienzo recibe el clip ya armado.}
    \label{fig:paquetes_extension}
\end{figure}

\subsection{Intercepción de cámara y lienzo virtual}
El script \texttt{inject-gum.js} se inyecta en el contexto principal de la página para sobrescribir \texttt{navigator.mediaDevices.getUserMedia}. La señal de la cámara física se deriva a un elemento \texttt{<video>} oculto y se redibuja en un \texttt{<canvas>} virtual a 24 cuadros por segundo. Mediante \texttt{canvas.captureStream()}, la extensión entrega este lienzo sintético a Google Meet, permitiendo superponer subtítulos sobre la franja inferior de la imagen transmitida por WebRTC.

\subsection{Motor local ILSA (FastAPI)}
El motor local se implementó en Python con FastAPI y se ejecuta sobre Uvicorn en la interfaz de bucle local (\texttt{127.0.0.1:8765}). El servidor configura reglas CORS específicas para permitir peticiones originadas únicamente en \texttt{https://meet.google.com} y la extensión local, incluyendo la cabecera \texttt{Access-Control-Allow-Private-Network: true}. La síntesis de voz se desacopló en un hilo secundario asíncrono para prevenir bloqueos en la atención de solicitudes HTTP.

\section{Implementación del avatar}

\subsection{Despliegue del servicio de archivos}
El servicio de archivos de poses (\texttt{tpose-server}) se empaquetó en una imagen de contenedor basada en Python 3.10-slim y se desplegó en la plataforma Google Cloud Run (región \texttt{southamerica-east1}). La arquitectura física de validación se representa en la Figura~\ref{fig:despliegue_prototipo}.

\begin{figure}[H]
    \centering
    \begin{tikzpicture}[
        font=\footnotesize,
        box/.style={draw, rounded corners=2pt, align=center, inner sep=4pt, fill=white, text width=3.6cm},
        arr/.style={-{Stealth[length=2.2mm]}, thick}
    ]
        \node[box] (chrome) {Chrome\\extensión de Meet};
        \node[box, below=0.6cm of chrome] (ilsa) {ILSA\\\texttt{127.0.0.1:8765}};
        \node[draw, rounded corners=3pt, inner sep=8pt, fit=(chrome)(ilsa),
              label={[font=\scriptsize]above:Equipo del usuario}] (pc) {};
        \node[box, right=2.4cm of chrome] (sem) {Traductor\\una instancia Cloud Run};
        \node[box, below=0.6cm of sem] (pose) {\texttt{tpose-server}\\\texttt{southamerica-east1}\\puerto 8080};
        \node[draw, rounded corners=3pt, inner sep=8pt, fit=(sem)(pose),
              label={[font=\scriptsize]above:Google Cloud Run}] (cloud) {};
        \draw[arr] (chrome) -- node[right, font=\scriptsize] {HTTP local} (ilsa);
        \draw[arr] (ilsa.east) -- node[above, font=\scriptsize] {texto} (sem.west);
        \draw[arr] (ilsa.south east) -- node[below, font=\scriptsize] {\texttt{GET /tposes}} (pose.west);
    \end{tikzpicture}
    \caption{Despliegue del prototipo validado experimentalmente.}
    \label{fig:despliegue_prototipo}
\end{figure}

\subsection{Vía descartada: retarget hacia un avatar VRM}
Se evaluó inicialmente la animación de un modelo tridimensional VRM mediante librerías de conversión de coordenadas a ángulos de Euler (Kalidokit) y cinemática inversa sobre Three.js. Esta aproximación se descartó debido a la inestabilidad en el cálculo de rotaciones provocada por la pérdida de profundidad absoluta en la coordenada $z$ de la cámara monocular. La aproximación por renderizado directo de coordenadas en formato \texttt{.pose} eliminó este artefacto al proyectar las posiciones espacializadas sin requerir cinemática inversa.

\subsection{Parámetros de composición y renderizado}
El ensamblado cinemático ejecuta el suavizado atenuado mediante la ecuación de promedio exponencial:
\begin{equation}
p_t = \alpha\, x_t + (1 - \alpha)\, p_{t-1}
\label{eq:ema_impl}
\end{equation}
con un factor $\alpha = 0{,}50$. El lienzo sintético se configura en $720 \times 780$ píxeles y escala la distancia entre hombros a $250$ píxeles. La duración de los puentes sintéticos interpola entre 4 y 8 cuadros según la distancia euclidiana entre las muñecas de los núcleos consecutivos.

\section{Instancia de validación del prototipo mediante túnel seguro}
La validación del sistema integrado se completó enlazando el motor ILSA local con los microservicios semánticos y de poses alojados en la nube a través de un túnel cifrado (\textit{ngrok}), confirmando el funcionamiento en tiempo pseudo-real sobre videollamadas activas en la plataforma Google Meet.